\section{What do we know about the impacts of AI sycophancy?}
In order to use norm displacement as a dividing line between people-pleasing behaviors and sycophancy in LLMs, it is critical that researchers understand the normative effects of these behaviors: both the direct impacts on output quality and the downstream impacts in the short-, medium-, and long-term. We detail these impacts in the following paragraphs.

\subsection{Direct Impacts on Output Quality}
Most of the existing works on AI sycophancy have sought to quantify the direct impacts of apparent approval-seeking behavior on output quality. Some of these impacts are more trivial to measure, like changes in accuracy \cite{sharma2024towards}, calibration \cite{sicilia-etal-2025-accounting} and hallucinations \cite{sharma2024towards,rrv-etal-2024-chaos}. Other impacts may be more difficult to measure: for instance, contradicting existing evidence in scenarios that may be subjective or where the model has incomplete information \cite{atwell2025basil}, or encouraging clinically harmful patterns in mental health support scenarios \cite{han2026detecting}.

Some works also study these impacts in more specific contexts, domains, or applications.
\citet{wang2026ai} find that AI reviewers were vulnerable to reversing their positions when presented with user rebuttals, even if this reversal was unfounded and not supported by the given evidence. \citet{lahlou2026writing} found that AI-generated literature reviews contained biases reflecting the users' perceived viewpoints or intended directions. 


\subsection{Downstream Impacts}
Few existing papers have identified  downstream impacts of AI sycophancy, and those that do have primarily focused on short-term impacts. We describe these works below, distinguishing between two categories of sycophantic behaviors defined in \citet{ye2026countsaisycophancytaxonomy}: \emph{Position} and \emph{Person}, where Position sycophancy involves conforming to or validating a user's opinions or beliefs, and Person sycophancy involves validating the person themselves (including character traits and emotions). 

\paragraph{Person sycophancy}
While many works speculate on the impacts of Person sycophancy on humans' psychology and social well-being, only a few works quantify these directly. \cite{lee2026effects} study the impact of sycophancy on social well-being on the context of AI companionship, and find a negative relationship between sycophantic model responses and perceived social support. \cite{chandra2026sycophantic} model the connection between AI sycophancy and AI psychosis, and find that even ideal Bayesian reasoners can experience delusional spiraling as a result of chatbot conversations, and that sycophancy plays a role in this phenomenon. \cite{Batista2026ARA} study how sycophancy can distort users' reality and reinforce user biases, even when they contradict the existing evidence. \cite{rathje2025sycophantic} study the effects of sycophancy in the context of political conversations, and find evidence indicating that humans prefer sycophantic models, and that engaging with sycophantic models in these discussions increased attitude extremity and certainty, indicating the possibility that sycophantic AI systems may contribute to echo chamber formation. \citet{cheng2026sycophantic} find that sycophantic responses in advice-seeking conversations regarding interpersonal conflicts lead to users feeling less inclined to repair relationships and more certain that they are in the right. They also found that users prefer sycophantic responses overall, despite the evidence that they impaired the users' social judgments. 

All of these works provide evidence that Person sycophancy may have a profound impact on humans' psychological health and social well-being, but there are not yet any works that capture this impact beyond the short-term (within an isolated set of conversations). Given the mounting evidence of AI psychosis and psychological harm stemming from long-term interactions with AI chatbots \cite{TheyAske50:online}, it is crucial that researchers conduct more medium-to-long-term studies on the impact of AI sycophancy on users' mental wellbeing, particularly in naturalistic settings. There is also a critical need to study how AI sycophancy may impact developing brains, as features which may help adult users may harm youth. For instance, \citet{jiao2026ai} point out that emotional mirroring, a behavior rewarded by adult users, can promote emotional dependency in younger users and impair emotional development by reducing social friction necessary to facilitate emotional development and develop executive functions. 

\paragraph{Position sycophancy}
There is also a growing body of work demonstrating the results of Position sycophancy (where the user's position is validated) on task success in various domains.

A few works have studied the impact of Position sycophancy in decision-making settings. \cite{li2026does} found that opinion conformity reduced the likelihood of users changing their mind when opposing evidence was provided. \citet{sicilia-etal-2025-accounting} found that sycophancy can reduce the success of human-AI collaboration for uncalibrated users. \citet{kumar2025ai} found evidence that AI software engineering agents often did not push back on or debate human developers, and that when they were directly encouraged to do so, developers uncovered new insights about their codebases. 

A couple of works have also studied the impact of Position sycophancy in multi-agent systems. \citet{kasprova2026too} found that providing agents with sycophancy priors in collaborative multi-agent systems improved discussion accuracy by 10.5\%. \citet{Yao2025PeacemakerOT} find that, in the setting of multi-agent debate, sycophancy amplifies disagreement collapse before reaching a correct conclusion and results in lower accuracy than single-debate settings. 

These works represent instances in which researchers have isolated the negative impacts of Position sycophancy in task-specific settings. A much larger body of literature exists which identifies sycophancy in various domains (mental health, finance, healthcare and medicine, software, etc.) but significantly fewer works have studied its impact on human-AI collaboration and task success. Given the critical, high-stakes nature of many of these domains, such as law, healthcare, and education, there is a need to conduct more assessments regarding the impact of AI sycophancy on specific processes within these domains. Further, many of the above works have studied the impacts of sycophancy on the micro level, isolating specific instances and interactions in which sycophancy has an impact on task success. Upon identifying these impacts at a granular level, it is crucial that we understand the systematic and institutional effects of AI sycophancy: its impacts on the peer-review process, patient care, etc.

\section{The case for a unified construct}

Many researchers characterize AI sycophancy as an inherently negative property, despite the critical gaps in the field's understanding about the downstream effects of people-pleasing behavior in LLMs. This makes it difficult to determine which approval-seeking behaviors are associated with negative impacts, and should thus fall under the umbrella of ``AI sycophancy''. As a result, most researchers focus on studying opinion shifts, particularly for objective tasks, as accuracy changes, factual errors, and hallucinations are easy to identify. However, recent works have indicated that other, more under-studied sycophantic behaviors, including social sycophancy \cite{cheng2026sycophantic}, can have negative downstream effects on users.  %, but these behaviors are comparatively under-studied in current research.

Many works on AI sycophancy are also motivated by its potential long-term effects, but these effects are largely unstudied and unverified. Some critical failures of generative AI systems have been posited to result from sycophancy in AI, but it is typically unclear what, if any, sycophantic behaviors are connected to these impacts.
Without a clear understanding of the downstream effects of a diverse set of sycophantic or people-pleasing behaviors in AI systems, it is difficult for researchers to make informed decisions about which behaviors are most important to quantify and mitigate to prevent certain harms.

We propose a standardized, bottom-up approach to characterizing and evaluating sycophancy in AI systems, where the focus is on identifying and quantifying norm displacements. This aligns with the behaviorist framing of existing research in AI sycophancy and serves as a means to resolve disagreements among researchers about the types of behaviors that should be characterized as sycophantic. It also encourages more careful, data-driven benchmark development that measures a diversity of sycophantic behaviors associated with downstream harms. Further, it gives researchers a clearer intuition regarding the possible broader impacts of different sycophancy mitigation strategies, allowing them to deploy selected mitigation strategies to target specific harms.

This framework is also in alignment with broader calls in the NLP and machine learning communities to consider downstream impacts when designing model evaluations. \citet{paullada2021data} have called out a number of problems associated with dataset development and benchmarking in AI disciplines, including  \citet{eriksson2025can} recently pointed out a number of problems with current AI benchmarks, including misalignments between what benchmarks actually studied and what they were claimed to study, a lack of a clear definition of the phenomena that benchmarks were designed to measure, lack of considerations for downstream impacts, and a lack of clarity regarding how benchmark results should be used in practice. 

% \begin{enumerate}[noitemsep]
%     \item Focus research towards studying sycophantic behaviors that are associated with the most negative impacts
%     \item Develop benchmarks that are informed by a more complete understanding of the behaviors that are most harmful
%     \item Evaluate targeted mitigation strategies, and understand which mitigation strategies are best for addressing which downstream outcomes
% \end{enumerate}